% Source: Weinberg GR, physical PDF pp. 639--642
%         (printed pp. 613--616; both boundary pages are partial).
% Coverage: Section 16.2, Models with a Cosmological Constant; equations
%           (16.2.1)--(16.2.15), including all unnumbered displays.
% Figures/tables/footnotes: linked reference markers 5--13; no figures or
%                          tables.
% Status: source-reviewed and compile-clean.
% Uncertainties: curvature, cosmological-constant, and scale-factor conversions
%                are documented below; no source uncertainty.
% MODERNIZATION CHECK: At fixed slots the target R_{\mu\nu}, R, and
% G_{\mu\nu} are the negatives of Weinberg's source curvature objects.
% Consequently (16.2.1)--(16.2.2) are written in the target form
% G_{\mu\nu}+\Lambda g_{\mu\nu}=8\pi G T_{\mu\nu}. The source
% cosmological-constant symbol is written \Lambda, and
% R_{\mathrm{source}}(t) is written a(t).

\subsection{Models with a Cosmological Constant}\label{sec:16.2}

When Einstein formulated the general theory of relativity in 1916, the universe
was generally believed to be static. According to Eqs.~\eqref{eq:15.1.18} and
\eqref{eq:15.1.19}, the scale factor \(a(t)\) can only be constant if
\[
  \rho=-3p=\frac{3k}{8\pi G a^2}.
\]
However, this requires that either the energy density \(\rho\) or the pressure
\(p\) should be \emph{negative}. In order to avoid this unphysical result,
Einstein in 1917 modified his equations to
read\hyperref[ch16-ref-5]{[5]}
\begin{equation}
  R_{\mu\nu}
  -\frac12g_{\mu\nu}R
  +\Lambda g_{\mu\nu}
  =
  8\pi G T_{\mu\nu},
  \tag{16.2.1}\label{eq:16.2.1}
\end{equation}
where \(\Lambda\) is a new fundamental constant, the so-called
\emph{cosmological constant}.

We have already noted at the end of Section~\ref{sec:7.1} that
Eq.~\eqref{eq:16.2.1} is the most general modification of Einstein's equations
that preserves the feature that \(T_{\mu\nu}\) is set equal to a tensor that is
constructed from \(g_{\mu\nu}\) and its first and second derivatives, and is
linear in the second derivatives of \(g_{\mu\nu}\). However, for our present
purposes it is more convenient to move \(\Lambda g_{\mu\nu}\) to the
right-hand side of the equations, writing
\begin{equation}
  R_{\mu\nu}
  -\frac12g_{\mu\nu}R
  =
  8\pi G\widetilde T_{\mu\nu},
  \tag{16.2.2}\label{eq:16.2.2}
\end{equation}
where \(\widetilde T_{\mu\nu}\) is a modified energy-momentum tensor:
\begin{equation}
  \widetilde T_{\mu\nu}
  \equiv
  T_{\mu\nu}
  -\frac{\Lambda}{8\pi G}g_{\mu\nu}.
  \tag{16.2.3}\label{eq:16.2.3}
\end{equation}
If \(T_{\mu\nu}\) has the perfect-fluid form \eqref{eq:15.1.12}, then so does
\(\widetilde T_{\mu\nu}\):
\begin{equation}
  \widetilde T_{\mu\nu}
  =
  \widetilde p\,g_{\mu\nu}
  +(\widetilde p+\widetilde\rho)u_\mu u_\nu,
  \tag{16.2.4}\label{eq:16.2.4}
\end{equation}
with a modified density and pressure
\begin{equation}
  \widetilde p
  \equiv
  p-\frac{\Lambda}{8\pi G},
  \qquad
  \widetilde\rho
  \equiv
  \rho+\frac{\Lambda}{8\pi G}.
  \tag{16.2.5}\label{eq:16.2.5}
\end{equation}
All of the results obtained in Section~\ref{sec:15.1} still apply to theories
with a cosmological constant, provided that we replace the quantities \(\rho\)
and \(p\) with the modified density and pressure \eqref{eq:16.2.5}.

In particular, the conditions for a static universe now read
\begin{equation}
  \widetilde\rho
  =
  -3\widetilde p
  =
  \frac{3k}{8\pi G a^2}.
  \tag{16.2.6}\label{eq:16.2.6}
\end{equation}
For a ``dust''-filled universe with \(p=0\), this gives
\begin{equation}
  \frac{k}{a^2}=\Lambda
  \tag{16.2.7}\label{eq:16.2.7}
\end{equation}
and
\begin{equation}
  \rho=\frac{\Lambda}{4\pi G}.
  \tag{16.2.8}\label{eq:16.2.8}
\end{equation}
In order to have \(\rho\) positive, Eq.~\eqref{eq:16.2.8} requires that
\(\Lambda\) should be positive, and Eq.~\eqref{eq:16.2.7} then tells us that
\begin{equation}
  k=+1
  \tag{16.2.9}\label{eq:16.2.9}
\end{equation}
and
\begin{equation}
  a=\frac{1}{\sqrt{\Lambda}}.
  \tag{16.2.10}\label{eq:16.2.10}
\end{equation}
The static Einstein universe is therefore \emph{finite} (though of course
unbounded), with a positive curvature and a density that are fixed by the
fundamental constants \(\Lambda\) and \(G\).

Of course, the discovery during the 1920's of a systematic relation between
red shift and distance removed any interest in the static Einstein universe as
a realistic cosmological model. Nevertheless, the existence of a cosmological
constant remains a logical possibility, and cosmologists have thoroughly
explored the dynamics of expanding universes with a cosmological
constant.\hyperref[ch16-ref-6]{[6]} We shall restrict our attention here to
models with zero pressure, so that Eq.~\eqref{eq:15.1.21} gives \(\rho a^3\)
constant. It is convenient to express this constant in terms of the value it
would have in a static Einstein model:
\begin{equation}
  \rho a^3
  =
  \frac{\alpha}{4\pi G\sqrt{|\Lambda|}}.
  \tag{16.2.11}\label{eq:16.2.11}
\end{equation}
\begin{sloppypar}
The dynamical equation \eqref{eq:15.1.20}, with \(\rho\) replaced by the
modified density \(\widetilde\rho\) defined in Eq.~\eqref{eq:16.2.5}, now reads
\end{sloppypar}
\begin{equation}
  \dot a^2
  =
  \frac1a
  \left\{
    \frac{\Lambda a^3}{3}
    -ka
    +\frac{2\alpha}{3\sqrt{|\Lambda|}}
  \right\}.
  \tag{16.2.12}\label{eq:16.2.12}
\end{equation}
The qualitative behavior of \(a(t)\) depends on the pattern of zeros, maxima,
and minima of the cubic on the right-hand side. There are three special cases
of particular interest, associated with the names of de Sitter, Lema\^{\i}tre,
and Eddington and Lema\^{\i}tre.

In the \emph{de Sitter model},\hyperref[ch16-ref-7]{[7]} space is essentially
empty and flat, so that \(k=\alpha=0\), and \(\Lambda\) is positive.
Equation~\eqref{eq:16.2.12} then has the simple solution
\begin{equation}
  a\propto e^{Ht},
  \tag{16.2.13}\label{eq:16.2.13}
\end{equation}
\begin{equation}
  H=\left(\frac{\Lambda}{3}\right)^{1/2}.
  \tag{16.2.14}\label{eq:16.2.14}
\end{equation}
The metric here is the same as in the steady state model discussed in
Section~\ref{sec:14.8}, with the difference that instead of matter being
continuously created, there is no matter at all! As discussed in
Section~\ref{sec:13.3}, this metric has a ten-parameter group of isometries,
which is just the group of ``rotations'' in five dimensions that leave
invariant a diagonal matrix with elements \(+1,+1,+1,+1,-1\). This group is
therefore often called the \emph{de Sitter group}. Although the absence of
matter in the de Sitter model removes it from consideration as a serious model
of the universe, it should be noted that any model with \(\Lambda>0\) goes over
to a de Sitter model for \(a\to\infty\).

In what is known as the \emph{Lema\^{\i}tre model},\hyperref[ch16-ref-8]{[8]}
space is positively curved, \(\Lambda\) is positive, and more matter is present
than in a static Einstein model, so that \(k=+1\), and \(\alpha>1\). According
to Eq.~\eqref{eq:16.2.12}, the scale factor \(a\) starts expanding at \(t=0\)
like \(t^{2/3}\), but the expansion then slows down, reaching a minimum rate at
\(a=\alpha^{1/3}/\sqrt{\Lambda}\), after which it speeds up again, ultimately
approaching the de Sitter result \eqref{eq:16.2.13}. The most remarkable
feature of this model is the existence of a ``coasting period'' during which
\(a(t)\) remains close to the value
\(a=\alpha^{1/3}/\sqrt{\Lambda}\) at which \(\dot a\) has its minimum. During
this period, the differential equation \eqref{eq:16.2.12} with \(k=+1\) takes
the approximate form
\[
  \dot a^2
  \simeq
  \alpha^{2/3}-1
  +\left(\sqrt{\Lambda}\,a-\alpha^{1/3}\right)^2.
\]
The solution is
\[
  a
  =
  \frac{\alpha^{1/3}}{\sqrt{\Lambda}}
  \left[
    1
    +(1-\alpha^{-2/3})^{1/2}
    \sinh\!\left(\sqrt{\Lambda}(t-t_m)\right)
  \right],
\]
where \(t_m\) is the time at which \(\dot a\) reaches its minimum. If \(\alpha\)
is very close to unity, then \(a\) will remain close to the static Einstein
value for a long time, of order
\begin{equation}
  \Delta t
  =
  \Lambda^{-1/2}
  \left|\ln(1-\alpha^{-2/3})\right|.
  \tag{16.2.15}\label{eq:16.2.15}
\end{equation}
The \emph{Eddington--Lema\^{\i}tre model} is a limiting case of the
Lema\^{\i}tre models, given particular prominence through the work of
Eddington.\hyperref[ch16-ref-9]{[9]} It has the same curvature and mass as the
static Einstein model; that is, \(k=+1\) and \(\alpha=1\), and behaves like a
Lema\^{\i}tre model with an infinitely long ``coasting period.''

Thus, if we start with \(a=0\) at \(t=0\), then \(a\) asymptotically approaches
the Einstein value \(1/\sqrt{\Lambda}\) for \(t\to\infty\). On the other hand,
if we start with \(a=1/\sqrt{\Lambda}\) for \(t=0\), then \(a\) grows
monotonically, ultimately approaching the de Sitter exponential growth
\eqref{eq:16.2.13}. This incidentally shows that the Einstein model is
\emph{unstable}, because if it is subjected to an infinitesimal expansion or
contraction, then \(a\) must go on expanding or contracting, with a time
dependence given by the Eddington--Lema\^{\i}tre model.

The observed concentration of the red shifts of quasi-stellar objects around
\(z\simeq2\) (see Sections~\ref{sec:11.6} and \ref{sec:14.8}) has revived
interest in the Lema\^{\i}tre models,\hyperref[ch16-ref-10]{[10]} since it
suggests that an unusually large number of QSO's were present at a particular
value of the scale factor, \(a\simeq a_0/3\), as would be expected in a
Lema\^{\i}tre model for which the ``coasting'' radius
\(\alpha^{1/3}/\sqrt{\Lambda}\) had this particular value. By taking \(\alpha\)
close to unity, we can make the ``coasting period'' as long as we want, so that
the predominance of the particular red shift \(z\simeq2\) can be made as
pronounced as may be needed to account for the QSO observations. With this new
motivation, there have recently been carried out studies of the propagation of
light signals around the universe,\hyperref[ch16-ref-11]{[11]} of the radio
number counts,\hyperref[ch16-ref-12]{[12]} and of the formation of
galaxies,\hyperref[ch16-ref-13]{[13]} in Lema\^{\i}tre models. There is no
definite evidence against the Lema\^{\i}tre models, but they seem an artificial
way to account for what may be merely a detailed feature of the evolution of
the quasi-stellar objects.
